% Generated draft: requires mathematical and source review.
\documentclass[11pt]{article}
\usepackage[margin=25mm]{geometry}
\usepackage{fontspec}
\setmainfont{Latin Modern Roman}
\usepackage{amsmath,amssymb,mathtools,unicode-math}
\setmathfont{Latin Modern Math}
\usepackage{enumitem,needspace,xurl,hyperref}
\setlength{\parindent}{0pt}
\setlength{\parskip}{5pt}
\setlength{\emergencystretch}{3em}
\newcommand{\N}{\mathbb{N}}
\newcommand{\Z}{\mathbb{Z}}
\newcommand{\Q}{\mathbb{Q}}
\newcommand{\R}{\mathbb{R}}
\newcommand{\C}{\mathbb{C}}
\newcommand{\F}{\mathbb{F}}
\newcommand{\A}{\mathbb{A}}
\newcommand{\PP}{\mathbb{P}}
\newcommand{\E}{\mathbb{E}}
\newcommand{\eps}{\varepsilon}
\newcommand{\dd}{\,\mathrm{d}}
\begin{document}
\section*{UnsolvedMath ID: 20002692}
\textbf{Source problem code:} catalogue entry
\subsection{Definitions and mathematical statement}
Question 32 asks for an explicit derivation $\partial$ whose positive generator $\partial^*\partial$ is the number operator of a $q$-deformed semicircular system, whose abstract existence the source attributes to Sauvageot. The source omits the $q$ range, target bimodule and operator domain. The aid supplies the following copied-letter realization for $-1<q<1$.

Let $H_{\R}$ be a real Hilbert space, $M=\Gamma_q(H_{\R})$ with its vacuum trace, and let $H_{\R}'$ be an orthogonal copy. Embed $M$ in
\[
\widetilde M=\Gamma_q(H_{\R}\oplus H_{\R}').
\]
Let $\alpha_t$ be second quantization of the orthogonal rotation sending $h$ to $\cos(t)h+\sin(t)h'$. On the Wick-polynomial algebra define $\partial x=(\dd/\dd t)_{t=0}\alpha_t(x)$, with values in the $M$-bimodule $L^2(\widetilde M)$ under the usual left and right actions. Explicitly,
\[
\partial W(h_1\otimes\cdots\otimes h_n)
=\sum_{j=1}^nW(h_1\otimes\cdots\otimes h_{j-1}\otimes h_j'\otimes h_{j+1}\otimes\cdots\otimes h_n).
\]
It satisfies $\partial(xy)=(\partial x)y+x(\partial y)$. The target can be restricted to the closed bimodule generated by Wick words with exactly one copied letter.

For the number operator $N$ equal to $n$ on the $n$th chaos, the aid's $q$-Fock pairing calculation gives
\[
\langle\partial x,\partial y\rangle_{L^2(\widetilde M)}
=\langle x,Ny\rangle_{L^2(M)}
\]
on Wick polynomials. The derivation is closable, its closure has domain $\operatorname{Dom}(N^{1/2})$, and $\overline\partial^{\,*}\overline\partial=N$. Equivalently, for a chaos expansion $x=\sum_{n\geq0}x_n$, the gradient domain is characterized by $\sum_n n\|x_n\|_2^2<\infty$. This fixes the domain and normalization, not merely a formal square-root identity.

In one variable let $X=s_q(e)$ for a unit vector $e$, $e'$ its copy, and use monic $q$-Hermite polynomials $H_0=1$, $H_1(x)=x$,
\[
xH_n(x)=H_{n+1}(x)+[n]_qH_{n-1}(x),\qquad [n]_q=\frac{1-q^n}{1-q}.
\]
Then $H_n(X)=W(e^{\otimes n})$ and
\[
\partial H_n(X)=\sum_{j=1}^nW(e^{\otimes(j-1)}\otimes e'\otimes e^{\otimes(n-j)}).
\]
At $q=0$ the copied-letter bimodule identifies with the usual direct sum of coarse bimodules and this derivation with the vector of free difference quotients.

The aid distinguishes this derivation from Fock annihilation. The lowering rule $aH_n=[n]_qH_{n-1}$ has gradient-energy ratio
\[
\frac{\|aH_n\|_2^2}{\|\partial H_n\|_2^2}=\frac{[n]_q}{n}.
\]
No fixed scalar multiple makes that ratio one for every $n$ when $-1<q<1$. In the ordinary multiplication bimodule it also violates Leibniz: $H_2(X)=X^2-I$, so Leibniz with $aX=I$ would give $2X$, whereas lowering gives $(1+q)X$. The formal exception $q=1$ lies outside the stated range and is not an added endpoint theorem.
\subsection{Short English statement}
The aid realizes the number-operator gradient by rotating into an orthogonal copied Gaussian system and differentiating. On a Wick word, replace each letter once by its copy and sum. This gives the correct energy and closed domain, while the familiar $q$-Hermite lowering operator has the wrong energy and Leibniz rule.
\subsection{Sources}
\begin{itemize}
\item Local UnsolvedMath registry, permanent ID $20002692$, \texttt{catalogue entry}; \texttt{aim-probability-notes.json}, \texttt{notes[133]}.
\item American Institute of Mathematics, \emph{Free Analysis}, workshop held 19--23 June 2006, Question 32, Section 0.11 ``Focus Group on Dirichlet Forms, from Classical to Quantum (day 5),'' according to supplied verification: \url{https://aimath.org/WWN/freeanalysis/freeanalysis.pdf}.
\item Sauvageot is the source attribution for abstract existence of the derivation. No title, theorem number or authoritative full citation is supplied for that construction or the copied-Fock formula.
\item Research attempt 1: full solution, problem solved at run, rigor label ``proved,'' novelty confidence low. Registry classification solved in literature, status solved, statement status exact. These labels remain supplied assessments requiring mathematical review.
\end{itemize}
\end{document}
